\documentclass[a4paper,landscape, 12pt]{extarticle}
\usepackage[left=0.5in, right=0.5in, top=0.4in, bottom=0.55in]{geometry}
\usepackage{graphicx} % Required for inserting images
\usepackage[normalem]{ulem}
\usepackage{tabto}
\usepackage{minted}
\renewcommand\ULthickness{1.0pt}   %%---> For changing thickness of underline
\setlength\ULdepth{1.3ex}%\maxdimen ---> For changing depth of underline
\usepackage{tikz}
\usetikzlibrary{graphs}
\usepackage{multicol}
\usepackage{array}
\usepackage{caption}
\usepackage{amsmath} 
\usepackage{fancyhdr}
\usepackage{lastpage}

\author{Juan R Ortiz}

% Page style with slide numbers
\pagestyle{fancy}
\fancyhf{}
\renewcommand{\headrulewidth}{0pt}
\fancyfoot[R]{\Large\thepage\ / \pageref{LastPage}}

% Define header command for consistent formatting
\newcommand{\slideheader}[2]{%
\noindent

\includegraphics[height=0.5in]{cpp_logo.png}%
\hspace{0.3in}%
{\Large \uline{#1 \hfill #2}}%
\par\vspace{0.4in}
}
% Standard content spacing - 0.6in base indent for all content
\newcommand{\contentstart}{%
\leftskip 0.6in \rightskip 0.6in
}
\newcommand{\contentend}{%
\leftskip 0in \rightskip 0in
}

% Standard font size for body content
\newcommand{\bodytext}{\LARGE}

\begin{document}
\begin{huge} 
%%%%%%%%%%%%%%%%%%%%%%%%%%%%%%%%%%%%%%%%
% SLIDE 1 - Title Slide
\slideheader{ECE5110 Numerical Modeling}{Unit 01}

\vspace{1.2in}
\begin{center}
{\Huge \textbf{Juan Ortiz}}

\vspace{0.4in}
{\LARGE Solving a function by the Newton Method}
\end{center}
%%%%%%%%%%%%%%%%%%%%%%%%%%%%%%%%%%%%%%%%
\newpage %Page 2
\slideheader{Newton Example: Designing a tone-shaping circuit for a guitar pedal}{Newton Method}

\leftskip 0.7in \rightskip 0.7in 
\vspace{0.4in}
You have been hired to help design the input stage of a new guitar effects pedal.
This pedal needs a low‑pass RC filter to remove harsh high‑frequencies while keeping the musical tone warm and natural.

\vspace{0.4in}
\noindent
The design team wants the filter’s cutoff frequency to be:
\vspace{0.1in}
\begin{center}$f_c = 1kHz$\end{center}
\vspace{0.1in}
The resistor value is already chosen:
\begin{center}$R = 10k\Omega$\end{center}
\vspace{0.1in}
Your job is to determine the capacitor value C that will produce the desired cutoff

\leftskip 0in \rightskip 0in 

%%%%%%%%%%%%%%%%%%%%%%%%%%%%%%%%%%%%%%%%
\newpage %Page 3
\slideheader{Newton Example: Designing a tone-shaping circuit for a guitar pedal}{Newton Method}

\leftskip 0.7in \rightskip 0.7in 
\vspace{0.4in}
\noindent
The cutoff frequency of a simple RC low‑pass filter is:
\begin{center}$f_c = \frac{1}{2\pi RC}$\end{center}
\vspace{0.2in}
This can be solved by rearranging the equation in root-finding form:
\begin{center}$g(C) = \frac{1}{2\pi RC} - f_c = 0$\end{center}
\vspace{0.2in}
Next define the function and it's derivative:
\vspace{0.2in}
\begin{center}$f(C) = \frac{1}{2\pi RC} - f_c$
\vspace{0.2in}
\\$f'(C) = -\frac{1}{2\pi RC^2}$\end{center}

\leftskip 0in \rightskip 0in 

%%%%%%%%%%%%%%%%%%%%%%%%%%%%%%%%%%%%%%%%
\newpage %Page 4
\slideheader{Newton Example: Designing a tone-shaping circuit for a guitar pedal}{Newton Method}

\leftskip 0.7in \rightskip 0.7in 
\vspace{0.4in}

In the guitar pedal world and analog audio circuits, the capacitance you choose in the signal path has a huge impact on the character of the sound.

\noindent
\\The standard capacitance for tone-shaping can range from 1nF-100nF, which will directly reflect electrical behavior and musical intent.
\begin{center}
\includegraphics[height=2in]{GTR_PDL.png}%
\end{center}
\begin{flushright}
1983 - BOSS DD-2
\end{flushright}

\leftskip 0in \rightskip 0in 

%%%%%%%%%%%%%%%%%%%%%%%%%%%%%%%%%%%%%%%%
\newpage %Page 5
\slideheader{Newton Example: Designing a tone-shaping circuit for a guitar pedal}{Newton Method}

\leftskip 0.7in \rightskip 0.7in 
\vspace{0.4in}
\noindent
Given the 1nF–100nF capacitance range used in guitar‑pedal signal paths, we can make a few reasonable assumptions when applying Newton’s Method to solve our filter equation.
\vspace{0.2in}
{
\begin{itemize}
    \item The plot interval: \{START, STOP\} = \{1nF, 100nF\} = \{1e-9, 1e-7\}
    \item Accuracy and tolerance: PRECISION = 1nF = 1e-9
    \item Initial guess: $x_0$ = 10nF = 1e-8
    \item Iterations: MAX\_STEPS = 20 (NM converges fast)
\end{itemize}
}
\vspace{0.2in} 
Recall the defined functions from the previous slides:
\vspace{0.2in}
{
\begin{itemize}
    \item f = 1/(2*math.pi*10000*C) - 1000
    \item df = -1/(2*math.pi*10000*C**2)
\end{itemize}
}
\leftskip 0in \rightskip 0in 

%%%%%%%%%%%%%%%%%%%%%%%%%%%%%%%%%%%%%%%%

\newpage %Page 6
\slideheader{Newton Example: Designing a tone-shaping circuit for a guitar pedal}{Newton Method}

\leftskip 0.7in \rightskip 0.7in 
\begin{minted}[fontsize=\LARGE]{python}
def solve_newton_method(self, f, df, x0, tolerance, max_steps):
    x = x0

    for i in range(1, max_steps + 1):
        dfx = df(x)
        if dfx == 0:
            return x, i, -1

        fx = f(x)
        x_new = x - fx / dfx

        if abs(x_new - x) < tolerance:
            return x_new, i, 0

        x = x_new

    return x, max_steps, -2
\end{minted}

\leftskip 0in \rightskip 0in 

%%%%%%%%%%%%%%%%%%%%%%%%%%%%%%%%%%%%%%%%
\newpage %Page 7
\slideheader{Newton Example: Designing a tone-shaping circuit for a guitar pedal}{Newton Method}

\leftskip 0.7in \rightskip 0.7in 
\vspace{0.4in}
Plug in the information VIA Python or Matlab:
\vspace{0.2in}
\begin{center}
\includegraphics[height=5.25in]{NumMet_Newton_GTRPDL.png}%
\end{center}
\leftskip 0in \rightskip 0in 

%%%%%%%%%%%%%%%%%%%%%%%%%%%%%%%%%%%%%%%%
\newpage %Page 8

% --- Background image using TikZ overlay (works in ANY document class) ---
\begin{tikzpicture}[remember picture,overlay]
    \node[opacity=0.15, at=(current page.center)] {
        \includegraphics[width=\paperwidth,height=\paperheight,keepaspectratio]{TS808.jpg}
    };
\end{tikzpicture}

% --- Your slide header (custom macro) ---
\slideheader{Newton Example: Designing a tone-shaping circuit for a guitar pedal}{Newton Method}

% --- Your slide content ---
\leftskip 0.7in 
\rightskip 0.7in

\begin{frame}{}

Our result is approximately a $16\,\text{nF}$ capacitor, which satisfies our low‑pass filter.

\vspace{0.2in}

This mid‑sized capacitor sits in a sweet spot for tone‑shaping, providing a gentle high‑frequency roll‑off used in classic pedals like the Big Muff, Tube Screamer, and Swollen Pickle.
\vspace{0.75in}
\begin{center}

\includegraphics[height=1.5in]{ECE_CAP.png}%
\end{center}

\end{frame}

\leftskip 0in 
\rightskip 0in

%%%%%%%%%%%%%%%%%%%%%%%%%%%%%%%%%%%%%%%%
\newpage %Page 9
\slideheader{Newton Example: Designing a tone-shaping circuit for a guitar pedal}{Newton Method}

\leftskip 0.7in \rightskip 0.7in 
\vspace{2.5in}
\begin{center}
{\Huge \textbf{END}}
\end{center}
\leftskip 0in \rightskip 0in 

%%%%%%%%%%%%%%%%%%%%%%%%%%%%%%%%%%%%%%%%


\end{huge}
\end{document}